	\begin{abstract}
		Si studiano procedure non adattive su \(N=b^m\) posizioni, con
		\(b,m\ge2\), costituite da \(m\) stadi di distribuzione ciclica in
		\(b\) mazzetti e raccolta secondo permutazioni assegnate. La scrittura in
		base \(b\), con cifre in \(D_b=\{0,\ldots,b-1\}\), e la rappresentazione
		mediante matrici di permutazione descrivono lo stadio come rotazione dei
		fattori tensoriali seguita da un'azione locale. Dopo un ciclo completo,
		gli esiti costituiscono la classe digitalmente separabile, di cardinalità
		\((b!)^m\), e ammettono una fattorizzazione di Kronecker unica in matrici
		di permutazione, fissati l'ordine dei fattori e la loro dimensione
		\(b\times b\).

		Le somme delle immagini sulle fibre digitali ricostruiscono i fattori locali
		e forniscono un criterio di riconoscimento fra tutte le permutazioni delle
		\(b^m\) posizioni: dati compatibili con un esito separabile ne determinano
		l'unicità anche nell'intero \(S_{b^m}\). Il rovesciamento globale e
		un'estensione algebrica con controlli locali indipendenti conservano la
		classe finale e introducono molteplicità uniformi, determinate
		esplicitamente. Il caso \(27=3^3\) illustra la teoria; un'appendice
		confronta il modello con le distribuzioni rettangolari e con la dinamica
		informativa del gioco delle ventuno carte.
	\end{abstract}

	\section{Introduzione}

	Il gioco delle ventisette carte appartiene alla famiglia dei procedimenti
	classici nei quali una distribuzione ripetuta traduce informazioni operative
	in cifre di una scrittura posizionale. L'identità
	\[
	27=3^3
	\]
	permette di distribuire il mazzo in tre mazzetti e di ripetere il gesto per
	tre stadi; le posizioni occupate, nelle raccolte successive, dai mazzetti che
	contengono una carta determinano progressivamente le cifre ternarie della sua
	posizione finale. Il legame con la notazione posizionale e con gli algoritmi
	di ordinamento per cifre è classico nella letteratura sul trucco di Gergonne
	\cite{bolker}.
	Analisi dinamiche di giochi affini e loro generalizzazioni compaiono in
	Champanerkar e Jani \cite{champanerkar-jani} e Quintero \cite{quintero}.

	Anche gli strumenti algebrici utilizzati qui appartengono a un quadro noto.
	Per i prodotti di Kronecker e le matrici di commutazione si rinvia a
	Henderson e Searle \cite{henderson-searle}; D'Angeli e Donno
	\cite{dangeli-donno} studiano le matrici di shuffle e il riordinamento dei
	fattori di prodotti di Kronecker. Le azioni coordinate su prodotti cartesiani
	si collocano nella teoria dei prodotti intrecciati in azione prodotto
	(\emph{wreath products in product action}) \cite{praeger-schneider}.
	Nel caso binario, Diaconis, Graham e Kantor
	\cite[Lemma~4]{diaconis-graham-kantor} descrivono esplicitamente, su
	\(2^m\) carte, un gruppo generato dai perfect shuffles isomorfo a
	\((C_2)^m\rtimes C_m\). Quel gruppo comprende anche la rotazione delle
	coordinate: va distinto dalla sola classe separabile degli esiti finali
	studiata qui, e le convenzioni operative vanno confrontate separatamente.

	Il modello qui studiato isola una versione non adattiva del procedimento. Si
	fissano interi
	\[
	b\ge 2,
	\qquad
	m\ge 2,
	\]
	e un mazzo di \(N=b^m\) carte. In ciascuno dei \(m\) stadi il mazzo viene
	distribuito ciclicamente in \(b\) mazzetti, conservando l'ordine interno, e i
	mazzetti vengono raccolti secondo una permutazione di \(S_b\) stabilita per
	quello stadio. È inoltre ammesso, in ciascuno stadio, un rovesciamento globale
	prefissato del mazzo, applicato prima della distribuzione. Il modello non
	considera controlli adattivi dipendenti dalla posizione di una carta. Nel sottocaso senza rovesciamenti, la successione delle
	permutazioni di raccolta forma direttamente un codice locale
	\[
	\boldsymbol{\rho}
	=
	(\rho_{m-1},\dots,\rho_0)
	\in S_b^m.
	\]
	Con i rovesciamenti, la storia fisica comprende anche una parola binaria
	\(\boldsymbol{\varepsilon}\in\{0,1\}^m\); il ciclo completo assorbe tali
	rovesciamenti in un codice locale effettivo, ancora appartenente a
	\(S_b^m\), che determina la permutazione globale delle \(b^m\) posizioni.

	Il problema specifico è comprendere come la distribuzione e la raccolta
	fisiche si traducano nel formalismo tensoriale. La domanda centrale non
	riguarda soltanto il calcolo della permutazione finale.
	Occorre riconoscere, fra tutti gli elementi di \(S_{b^m}\), quelli che possono
	nascere dal procedimento, individuare la fattorizzazione locale nascosta nella
	loro azione e stabilire quali informazioni aggregate siano sufficienti a
	ricostruirla. La scrittura posizionale, il prodotto tensoriale e l'azione del
	prodotto diretto \(S_b^m\) convergono così nello stesso problema.
	Le somme di fibra sono considerate come dati aggregati della permutazione,
	senza presupporne proprietà di minimalità o vantaggi computazionali.

	Il ciclo fisico viene dapprima descritto mediante una rotazione dei fattori
	tensoriali; dopo \(m\) stadi la rotazione si chiude e la matrice globale si
	separa in un prodotto di Kronecker. Il rovesciamento globale, essendo a sua
	volta un prodotto di Kronecker con fattori tutti uguali, commuta con tale
	rotazione e non distrugge la separabilità digitale. Più in generale, fattori
	di rovesciamento diversi sui diversi livelli vengono trasportati ciclicamente
	dalla rotazione e restano decomponibili; anche il modello tensoriale esteso si
	chiude quindi nella medesima classe globale.

	La classe risultante ammette una caratterizzazione intrinseca: una
	permutazione globale vi appartiene esattamente quando preserva, livello per
	livello, le partizioni digitali delle posizioni. I fattori locali sono allora
	determinati in modo unico. Le somme delle immagini lungo le fibre digitali
	consentono inoltre di ricostruire tali fattori e forniscono un criterio di
	riconoscimento che non presuppone la separabilità della permutazione assegnata.

	I contributi specifici del lavoro seguono questa progressione. La sezione~2
	descrive il ciclo non adattivo mediante cifre; la sezione~3 rappresenta lo
	stadio come rotazione e azione locale e deduce la fattorizzazione dopo un
	ciclo completo. La sezione~4 caratterizza gli esiti digitalmente separabili;
	la sezione~5 analizza rovesciamenti, fattori effettivi e molteplicità delle
	storie operative, distinguendo il modello fisico dall'estensione algebrica.
	La sezione~6 sviluppa il problema inverso: formula delle somme di fibra,
	ricostruzione dei fattori e criterio di riconoscimento. In particolare, il
	Teorema~\ref{thm:criterio-fibra} mostra che le somme non consentono soltanto
	di ricostruire i fattori entro la classe separabile, ma caratterizzano ogni
	suo elemento anche fra tutte le permutazioni delle \(b^m\) posizioni.
	Il collegamento sistematico fra procedimento fisico, classe coordinata e
	criterio inverso costituisce il contributo del lavoro; gli strumenti
	tensoriali e di teoria dei gruppi restano quelli standard richiamati sopra.

	\section{Modello posizionale e procedura ordinaria}

	Iniziamo dal sottocaso ordinario, nel quale non intervengono rovesciamenti.
	Questa scelta serve a isolare la ricorrenza posizionale di base; la sezione
	dedicata ai rovesciamenti mostrerà che il rovesciamento globale può essere aggiunto senza
	uscire dalla classe delle trasformazioni digitalmente separabili.

	Sia
	\[
	\Db=\{0,1,\dots,b-1\}.
	\]
	Le posizioni del mazzo sono numerate da \(0\) a \(b^m-1\). Ogni posizione
	possiede una rappresentazione unica
	\[
	x=(a_{m-1}a_{m-2}\cdots a_1a_0)_b,
	\qquad
	a_i\in\Db,
	\]
	ossia
	\[
	x=a_{m-1}b^{m-1}+a_{m-2}b^{m-2}+\cdots+a_1b+a_0.
	\]
	Per
	\[
	\alpha=(a_{m-1},\dots,a_0)\in\Db^m
	\]
	si pone
	\[
	\val(\alpha)
	=
	\sum_{i=0}^{m-1}a_i b^i.
	\]
	La funzione \(\val\) è una biiezione fra \(\Db^m\) e
	\(\{0,\dots,b^m-1\}\).

	Durante un singolo stadio, la posizione
	\[
	x=bq+r,
	\qquad
	r=x\bmod b,
	\qquad
	q=\left\lfloor\frac{x}{b}\right\rfloor
	\]
	viene distribuita nel mazzetto identificato dal resto \(r\). Se una
	permutazione \(\rho\in S_b\) colloca quel mazzetto nella posizione di blocco
	\(\rho(r)\), la nuova posizione è
	\[
	x'
	=
	\rho(r)b^{m-1}
	+
	\left\lfloor\frac{x}{b}\right\rfloor.
	\]
	La cifra meno significativa viene così estratta dall'indirizzo e reinserita
	come cifra più significativa, dopo l'azione della permutazione di raccolta.

	\begin{proposition}[Ricorrenza posizionale]
		Siano \(\rho_0,\dots,\rho_{k-1}\in S_b\) le permutazioni usate nei primi
		\(k\) stadi, con \(1\le k\le m\), e sia
		\[
		x_0=(a_{m-1}\cdots a_0)_b.
		\]
		Allora
		\[
		\begin{aligned}
			x_k
			={}&
			\rho_{k-1}(a_{k-1})b^{m-1}
			+\rho_{k-2}(a_{k-2})b^{m-2}
			+\cdots
			+\rho_0(a_0)b^{m-k}
			+\left\lfloor\frac{x_0}{b^k}\right\rfloor
		\end{aligned}
		\]
		In particolare, dopo \(m\) stadi,
		\[
		x_m
		=
		\sum_{i=0}^{m-1}\rho_i(a_i)b^i.
		\]
	\end{proposition}

	\begin{proof}
		Per \(k=1\) la formula coincide con la regola del singolo stadio. Se essa
		vale per un certo \(k<m\), la cifra successivamente estratta dal termine
		\(\lfloor x_0/b^k\rfloor\) è \(a_k\), mentre ciascuno dei termini già
		prodotti è divisibile per \(b\). La relazione
		\[
		x_{k+1}
		=
		\rho_k(a_k)b^{m-1}
		+
		\left\lfloor\frac{x_k}{b}\right\rfloor
		\]
		fornisce quindi la formula con \(k+1\) al posto di \(k\). Per \(k=m\), il
		termine \(\lfloor x_0/b^m\rfloor\) è nullo e le cifre trasformate occupano i
		livelli originari.
	\end{proof}

	Nel caso \(b=m=3\), la formula descrive il ciclo completo del gioco delle
	ventisette carte. Il suo significato tensoriale permette di distinguere con
	precisione l'operatore di uno stadio dalla trasformazione finale.

	\section{Dinamica tensoriale del ciclo completo}

	Sia
	\[
	E_b=\mathbb R^b
	\]
	con base canonica \(e_0,\dots,e_{b-1}\). All'indirizzo
	\(\alpha=(a_{m-1},\dots,a_0)\) si associa il tensore
	\[
	e_\alpha
	=
	e_{a_{m-1}}\otimes\cdots\otimes e_{a_0}
	\in E_b^{\otimes m}.
	\]
	La distribuzione ciclica è rappresentata dall'operatore
	\[
	\MSC_m:E_b^{\otimes m}\longrightarrow E_b^{\otimes m}
	\]
	definito sui tensori di base da
	\[
	\MSC_m
	\left(
	e_{a_{m-1}}\otimes\cdots\otimes e_{a_0}
	\right)
	=
	e_{a_0}\otimes e_{a_{m-1}}\otimes\cdots\otimes e_{a_1}.
	\]
	Esso realizza una rotazione ciclica dei fattori e soddisfa
	\[
	\MSC_m^m=I.
	\]

	Identifichiamo la base tensoriale con la base delle posizioni mediante

	\[
	 e_{a_{m-1}}\otimes\cdots\otimes e_{a_0}
	 \longleftrightarrow
	 e_{\val(a_{m-1},\dots,a_0)},
	\]

	ordinando quindi i vettori di base secondo il valore crescente di \(\val\). Con questa identificazione, la matrice tensoriale rappresenta direttamente l'azione sulle posizioni numerate da \(0\) a \(b^m-1\).

	Per ogni \(\sigma\in S_b\), indichiamo con \(P_\sigma\) la matrice di
	permutazione definita da
	\[
	P_\sigma e_r=e_{\sigma(r)}.
	\]
	Quando non vi è ambiguità, identificheremo una permutazione con la relativa
	matrice. In particolare, poniamo \(S_h=P_{\rho_h}\).
	L'operatore dello stadio \(h\), con \(h=0,\dots,m-1\), è
	\[
	T_h
	=
	(S_h\otimes I_b^{\otimes(m-1)})\circ\MSC_m.
	\]
	La rotazione porta la cifra da elaborare nel primo fattore; la raccolta
	applica poi \(S_h\) a quel fattore. La separazione delle azioni locali emerge
	soltanto dopo il ciclo completo.

	\begin{theorem}[Chiusura tensoriale del ciclo]
		Con le convenzioni precedenti,
		\[
		T_{m-1}\circ\cdots\circ T_1\circ T_0
		=
		S_{m-1}\otimes S_{m-2}\otimes\cdots\otimes S_0.
		\]
	\end{theorem}

	\begin{proof}
		Applicando successivamente gli operatori a
		\[
		e_{a_{m-1}}\otimes\cdots\otimes e_{a_0},
		\]
		dopo il primo stadio si ottiene
		\[
		e_{\rho_0(a_0)}\otimes e_{a_{m-1}}\otimes\cdots\otimes e_{a_1},
		\]
		e dopo il secondo
		\[
		e_{\rho_1(a_1)}\otimes e_{\rho_0(a_0)}
		\otimes e_{a_{m-1}}\otimes\cdots\otimes e_{a_2}.
		\]
		Proseguendo per \(m\) stadi, ciascuna cifra viene esposta una volta
		all'azione della permutazione corrispondente e i fattori ritornano
		nell'ordine iniziale. Il risultato è
		\[
		e_{\rho_{m-1}(a_{m-1})}
		\otimes\cdots\otimes
		e_{\rho_0(a_0)},
		\]
		che coincide con l'azione di
		\(S_{m-1}\otimes\cdots\otimes S_0\). I tensori elementari formano una base,
		perciò i due operatori sono uguali.
	\end{proof}

	La relazione \(\MSC_m^m=I\) non autorizza una cancellazione formale delle
	rotazioni nella composizione, poiché fra due rotazioni successive compaiono
	le trasformazioni locali. Il teorema descrive invece il trasporto di ogni
	fattore fino al livello sul quale deve agire la raccolta e la successiva
	chiusura dell'intero ciclo.
	La rappresentazione dello stadio permette dunque di seguire la catena
	distribuzione e raccolta, rotazione, azione locale, trasporto e chiusura nel
	prodotto di Kronecker finale. Il Teorema~3.1 applica al procedimento
	considerato il meccanismo noto di riordinamento dei fattori mediante matrici
	di shuffle \cite{dangeli-donno}.

	\section{Permutazioni digitalmente separabili}

	D'ora in poi \(\boldsymbol{\rho}\) denota il codice locale di una
	trasformazione digitalmente separabile:
	\[
	\boldsymbol{\rho}
	=
	(\rho_{m-1},\dots,\rho_0)
	\in S_b^m.
	\]
	Esso definisce la trasformazione globale
	\[
	\pi_{\boldsymbol{\rho}}
	\bigl(\val(a_{m-1},\dots,a_0)\bigr)
	=
	\val\bigl(
	\rho_{m-1}(a_{m-1}),\dots,\rho_0(a_0)
	\bigr).
	\]
	La matrice di permutazione corrispondente, nella base ordinata dai tensori
	\(e_\alpha\), è
	\[
	P_{\boldsymbol{\rho}}
	=
	Q_{m-1}\otimes\cdots\otimes Q_0,
	\qquad Q_i=P_{\rho_i}.
	\]

	Per ogni livello \(i\in\{0,\dots,m-1\}\), sia
	\[
	d_i\bigl(\val(a_{m-1},\dots,a_0)\bigr)=a_i
	\]
	la proiezione sulla cifra di livello \(i\). Per \(t\in\Db\) si definisce la
	fibra digitale
	\[
	F_{i,t}
	=
	\{x\in\{0,\dots,b^m-1\}:d_i(x)=t\},
	\]
	e la partizione
	\[
	\mathcal F_i
	=
	\{F_{i,t}:t\in\Db\}.
	\]
	Ogni fibra ha cardinalità \(b^{m-1}\).

	\begin{definition}
		Una permutazione \(\pi\in S_{b^m}\) è \emph{digitalmente separabile} se
		esistono permutazioni \(\rho_0,\dots,\rho_{m-1}\in S_b\) tali che
		\[
		d_i(\pi(x))
		=
		\rho_i(d_i(x))
		\]
		per ogni posizione \(x\) e ogni livello \(i\).
	\end{definition}

	La proprietà può essere riconosciuta direttamente nell'azione di \(\pi\),
	senza conoscere in anticipo il codice che l'ha prodotta.

	\begin{theorem}[Caratterizzazione intrinseca e fattorizzazione]\label{thm:caratterizzazione-separabile}
		Per una permutazione \(\pi\in S_{b^m}\), le seguenti proprietà sono
		equivalenti.
		\begin{enumerate}
			\item Esiste un codice unico
			\[
			\boldsymbol{\rho}
			=(\rho_{m-1},\dots,\rho_0)\in S_b^m
			\]
			tale che \(\pi=\pi_{\boldsymbol{\rho}}\).

			\item Per ogni \(i\), la funzione \(d_i\circ\pi\) è costante su ciascuna
			fibra \(F_{i,t}\).

			\item Per ogni \(i\), la permutazione \(\pi\) induce una permutazione dei
			blocchi della partizione \(\mathcal F_i\).

			\item La matrice di permutazione di \(\pi\) ammette una fattorizzazione
			\[
			Q_{m-1}\otimes\cdots\otimes Q_0,
			\]
			dove ciascun \(Q_i\) è una matrice di permutazione \(b\times b\).
		\end{enumerate}
		Le permutazioni \(\rho_i\) e le matrici \(Q_i\) sono determinate in modo
		unico.
	\end{theorem}

	\begin{proof}
		La prima proprietà implica la seconda perché, per
		\(x=\val(a_{m-1},\dots,a_0)\), vale
		\[
		d_i(\pi(x))=\rho_i(a_i).
		\]
		Il membro destro dipende soltanto dal valore \(a_i=d_i(x)\), perciò è
		costante su \(F_{i,t}\).

		Supposta la seconda proprietà, per ogni \(i\) e \(t\) si definisce
		\[
		\rho_i(t)
		=
		d_i(\pi(x)),
		\qquad x\in F_{i,t}.
		\]
		La definizione non dipende dalla scelta di \(x\). Se due valori distinti
		\(t,t'\) avessero la stessa immagine \(u\), l'iniettività di \(\pi\) dovrebbe
		mandare i due insiemi disgiunti \(F_{i,t}\) e \(F_{i,t'}\), ciascuno di
		cardinalità \(b^{m-1}\), dentro la sola fibra \(F_{i,u}\), che possiede la
		stessa cardinalità di uno di essi. Ciò è impossibile. La funzione \(\rho_i\)
		è quindi iniettiva e, poiché \(\Db\) è finito, appartiene a \(S_b\).

		Per ogni posizione \(x\), le cifre di \(\pi(x)\) sono allora
		\[
		\rho_{m-1}(d_{m-1}(x)),\dots,\rho_0(d_0(x)).
		\]
		Esse determinano un'unica posizione, e ne segue
		\(\pi=\pi_{\boldsymbol{\rho}}\). La stessa costruzione mostra l'unicità di
		ciascuna \(\rho_i\).

		La seconda e la terza proprietà sono equivalenti: la costanza di
		\(d_i\circ\pi\) su \(F_{i,t}\) equivale all'inclusione dell'immagine della
		fibra in una singola fibra finale; la biiettività di \(\pi\) e l'uguaglianza
		delle cardinalità trasformano l'inclusione in uguaglianza.

		Infine, dalla prima proprietà e dall'azione sui tensori della base canonica
			segue direttamente
			\[
			P_{\boldsymbol{\rho}}
			\bigl(e_{a_{m-1}}\otimes\cdots\otimes e_{a_0}\bigr)
			=
			e_{\rho_{m-1}(a_{m-1})}\otimes\cdots\otimes e_{\rho_0(a_0)}.
			\]
			Pertanto
			\[
			P_{\boldsymbol{\rho}}
			=
			Q_{m-1}\otimes\cdots\otimes Q_0,
			\qquad Q_i=P_{\rho_i}.
			\]
			Viceversa, un prodotto
			\(Q_{m-1}\otimes\cdots\otimes Q_0\), con ciascun \(Q_i\) matrice
			di permutazione \(b\times b\), agisce separatamente sui rispettivi
			indici e induce una permutazione \(\pi_{\boldsymbol{\rho}}\).
			L'unicità delle permutazioni locali comporta quella dei fattori matriciali
			nell'ordine tensoriale fissato.
	\end{proof}

	\begin{corollary}
		L'applicazione
		\[
		\Phi:S_b^m\longrightarrow S_{b^m},
		\qquad
		\boldsymbol{\rho}\longmapsto\pi_{\boldsymbol{\rho}},
		\]
		è un omomorfismo iniettivo. La sua immagine coincide con il gruppo delle
		permutazioni digitalmente separabili ed è isomorfa a \(S_b^m\).
	\end{corollary}

	\begin{proof}
		La composizione di due codici avviene componente per componente, perciò
		\[
		\pi_{\boldsymbol{\rho}\boldsymbol{\sigma}}
		=
		\pi_{\boldsymbol{\rho}}\circ\pi_{\boldsymbol{\sigma}}.
		\]
		L'unicità del codice associato a una trasformazione rende il nucleo banale.
		La descrizione dell'immagine segue dal teorema.
	\end{proof}

	Poniamo
	\[
	\mathcal H_{b,m}=\{Q_{m-1}\otimes\cdots\otimes Q_0:Q_i\in S_b\}
	\cong S_b^m.
	\]
	La rappresentazione mediante fattori locali nell'ordine tensoriale fissato è
	iniettiva e
	\[
	|\mathcal H_{b,m}|=(b!)^m.
	\]
	Nella naturale azione su \(\Db^m\), questo è il gruppo base \(S_b^m\)
	del prodotto intrecciato \(S_b\wr S_m\) in azione prodotto
	\cite[sezione~1.1]{praeger-schneider}. Si tratta di una struttura standard
	di azioni coordinate. Dal Teorema~3.1 segue inoltre che, al variare delle
	\(m\) raccolte locali, il ciclo fisico realizza precisamente tutti gli
	elementi di questo gruppo base. La classe degli esiti finali non comprende le
	ulteriori permutazioni dei posti ammesse dal prodotto intrecciato completo.

	Nel caso \(b=m=3\), si ottiene
	\[
	S_3^3\hookrightarrow S_{27}.
	\]
	L'immagine contiene \(6^3=216\) permutazioni, una sottoclasse fortemente
	vincolata dei \(27!\) elementi di \(S_{27}\).

	
\section{Rovesciamenti e forma fondamentale dello stadio}

	La procedura con rovesciamenti viene descritta senza cambiare
	l'architettura dello stadio. Il punto di partenza, sia nel modello ordinario
	sia nella famiglia tensoriale estesa, è sempre
	\[
	\boxed{
	T_h
	=
	P_h\circ\MSC_m\circ J_h,
	}
	\qquad
	h=0,\dots,m-1.
	\]
	Il blocco \(J_h\) agisce per primo, quindi il rovesciamento precede la
	distribuzione; \(P_h\) rappresenta invece le permutazioni applicate dopo la
	distribuzione. Le diverse famiglie considerate in seguito non modificano
	questa formula: cambiano soltanto i vincoli imposti ai fattori locali di
	\(P_h\) e \(J_h\).

	\subsection{Rovesciamento locale e trasporto}

	Sul singolo livello digitale definiamo
	\[
	r_b(t)=b-1-t,
	\qquad
	t\in\Db,
	\]
	e indichiamo con \(R_b\) la corrispondente matrice di permutazione,
	\[
	R_be_t=e_{b-1-t}.
	\]
	Vale
	\[
	R_b^2=I_b.
	\]
	Il rovesciamento globale delle \(b^m\) posizioni è
	\[
	j_b(x)=b^m-1-x.
	\]
	Se
	\[
	x=(a_{m-1}\cdots a_0)_b,
	\]
	allora
	\[
	j_b(x)
	=
	\bigl(
	(b-1-a_{m-1})\cdots(b-1-a_0)
	\bigr)_b.
	\]
	Nella base tensoriale il corrispondente operatore è dunque
	\[
	J_b
	=
	R_b^{\otimes m}.
	\]

	Per un operatore decomponibile
	\[
	B
	=
	B_{m-1}\otimes\cdots\otimes B_0
	\]
	definiamo
	\[
	\vartheta(B)
	=
	B_0\otimes B_{m-1}\otimes\cdots\otimes B_1.
	\]
	Dalla definizione di \(\MSC_m\) segue
	\[
	\boxed{
	\MSC_m\circ B
	=
	\vartheta(B)\circ\MSC_m.
	}
	\]
	Equivalentemente,
	\[
	\vartheta(B)
	=
	\MSC_m\circ B\circ\MSC_m^{-1}.
	\]
	La coniugazione mediante \(\MSC_m\) definisce un automorfismo
	dell'algebra degli endomorfismi di \(E_b^{\otimes m}\). In particolare,
	\[
	\vartheta(BC)
	=
	\vartheta(B)\vartheta(C),
	\qquad
	\vartheta^m=\mathrm{id}.
	\]
	La sua restrizione a \(\mathcal H_{b,m}\) è un automorfismo di gruppo.


	\subsection{Il modello ordinario con rovesciamento globale}

	Nel modello ordinario la raccolta agisce su un solo livello alla volta.
	Indichiamo con \(S_h\in S_b\) la matrice della raccolta fisica allo stadio
	\(h\). I due blocchi della forma fondamentale sono
	\[
	P_h
	=
	S_h\otimes I_b^{\otimes(m-1)}
	\]
	e
	\[
	J_h
	=
	J_b^{\varepsilon_h}
	=
	(R_b^{\varepsilon_h})^{\otimes m},
	\qquad
	\varepsilon_h\in\{0,1\}.
	\]
	Lo stadio resta quindi
	\[
		T_h
	=
	(S_h\otimes I_b^{\otimes(m-1)})
	\circ\MSC_m\circ
	(R_b^{\varepsilon_h})^{\otimes m}.
	\]

	Poiché tutti i fattori di \(J_h\) coincidono,
	\[
	\vartheta(J_h)=J_h,
	\]
	e quindi
	\[
	\MSC_m\circ J_h
	=
	J_h\circ\MSC_m.
	\]
	Questa commutazione non permette di spostare \(J_h\) attraverso le raccolte
	\(P_k\); consente però di seguire esattamente la composizione degli stadi.

	\begin{proposition}[Molteplicità del modello ordinario con rovesciamento globale]
		\label{prop:molteplicita-ordinaria}
		Fissata una storia
		\[
		\boldsymbol{\varepsilon}
		=
		(\varepsilon_{m-1},\dots,\varepsilon_0)
		\in\{0,1\}^m,
		\]
		la mappa dalle \(m\) raccolte
		\[
		(S_{m-1},\dots,S_0)\in S_b^m
		\]
		alla trasformazione finale digitalmente separabile è biunivoca.
		Di conseguenza ogni trasformazione finale possiede esattamente
		\(2^m\) realizzazioni nel modello ordinario quando la storia dei
		rovesciamenti è considerata parte della procedura.
	\end{proposition}

	\begin{proof}[Dimostrazione per composizione degli stadi]
		Si parte da
		\[
		\begin{aligned}
		T
		&=
		T_{m-1}\circ\cdots\circ T_0\\
		&= (P_{m-1}\circ\MSC_m\circ J_{m-1})
		\circ\cdots\circ
		(P_0\circ\MSC_m\circ J_0).
		\end{aligned}
		\]
		Usando \(\MSC_m\circ J_h=J_h\circ\MSC_m\) in ogni stadio,
		\[
		T_h
		=
		P_h\circ J_h\circ\MSC_m.
		\]
		Trasportando successivamente attraverso le copie di \(\MSC_m\) si ottiene
		\[
		\begin{aligned}
		T
		={}&
		(P_{m-1}\circ J_{m-1})
		\circ
		\vartheta(P_{m-2}\circ J_{m-2})
		\circ \cdots \circ\vartheta^{m-1}(P_0\circ J_0)
		\end{aligned}
		\]
		perché \(\MSC_m^m=I\).

		Per ciascun \(h\),
		\[
				\vartheta^{m-1-h}(P_h)
		=
		I_b^{\otimes(m-1-h)}
		\otimes S_h\otimes
		I_b^{\otimes h}.
		\]
		Il fattore trasportato contiene quindi una sola permutazione variabile
		\(S_h\), collocata su un livello tensoriale diverso da quelli occupati
		dalle altre raccolte. Tutti i fattori provenienti dai \(J_h\) sono invece
		fissati dalla storia \(\boldsymbol{\varepsilon}\).
		Su ogni livello finale compare quindi esattamente una delle variabili
		\(S_h\), moltiplicata a sinistra e a destra per matrici fissate appartenenti
		a \(\{I_b,R_b\}\).

		La moltiplicazione a sinistra o a destra per una permutazione fissata è
		una biiezione di \(S_b\) in sé. Perciò, fissata
		\(\boldsymbol{\varepsilon}\), ciascuna \(S_h\) determina
		biunivocamente il fattore locale finale del livello che essa raggiunge.
		I \(m\) livelli sono indipendenti, dunque la mappa complessiva
		\(S_b^m\to S_b^m\) è biunivoca.

		Le storie binarie possibili sono \(2^m\). Per ognuna di esse esiste
		quindi una e una sola \(m\)-upla di raccolte che realizza una
		trasformazione finale assegnata, da cui la molteplicità \(2^m\).
	\end{proof}

	Nella formula seguente \(S_h\) indica la raccolta fisicamente scelta allo
	stadio \(h\), mentre \(Q_i\) è il fattore locale effettivo al livello
	\(i\) dell'esito finale, comprensivo dell'effetto dei rovesciamenti.

	\begin{corollary}[Formula esplicita e inversione delle raccolte]
		\label{cor:formula-inversione-raccolte}
		Fissata la storia \(\boldsymbol{\varepsilon}\), poniamo, con somme
		modulo \(2\),
		\[
		u_i
		=
		\sum_{h=0}^{i}\varepsilon_h,
		\qquad
		v_i
		=
		\sum_{h=i+1}^{m-1}\varepsilon_h.
		\]
		Se
		\[
		T=Q_{m-1}\otimes\cdots\otimes Q_0,
		\]
		allora, per ogni livello \(i\),
		\[
		\boxed{
		Q_i
		=
		R_b^{v_i}\circ S_i\circ R_b^{u_i}
		}
		\qquad\text{e}\qquad
		\boxed{
		S_i
		=
		R_b^{v_i}\circ Q_i\circ R_b^{u_i}.
		}
		\]
	\end{corollary}

	\begin{proof}
		La cifra originaria di livello \(i\) subisce, prima della raccolta
		\(S_i\), i rovesciamenti degli stadi \(0,\dots,i\), e dopo la raccolta
		quelli degli stadi \(i+1,\dots,m-1\). Il loro contributo complessivo è
		quindi rispettivamente \(R_b^{u_i}\) e \(R_b^{v_i}\), da cui
		\[
		Q_i
		=
		R_b^{v_i}\circ S_i\circ R_b^{u_i}.
		\]
		Poiché \(R_b^{-1}=R_b\), moltiplicando a sinistra e a destra per le
		stesse involuzioni si ricava la formula inversa per \(S_i\).
	\end{proof}

	\subsection{La famiglia tensoriale estesa}

	La famiglia dei controlli può essere ampliata mantenendo invariata
	l'architettura dello stadio. Usiamo il gruppo
	\(\mathcal H_{b,m}\cong S_b^m\) introdotto nella sezione precedente.
	Si tratta di un'estensione tensoriale e algebrica: i controlli indipendenti
	su tutti i fattori non rappresentano semplicemente le normali raccolte
	fisiche del gioco. Le formule e i conteggi seguenti riguardano lo spazio
	dei controlli qui definito, senza presupporne una realizzazione fisica.

	Nella famiglia estesa, allo stadio \(h\),
	\[
	P_h
	=
	P_{m-1,h}\otimes\cdots\otimes P_{0,h},
	\qquad
	P_{i,h}\in S_b,
	\]
	e
	\[
	J_h
	=
	R_b^{\varepsilon_{m-1,h}}
	\otimes\cdots\otimes
	R_b^{\varepsilon_{0,h}},
	\qquad
	\varepsilon_{i,h}\in\{0,1\}.
	\]
	Nei doppi indici, \(i\) identifica il posto tensoriale e \(h\) lo stadio;
	\(P_{i,h}\) è una matrice \(b\times b\), mentre il blocco \(P_h\) ha
	dimensione \(b^m\times b^m\).
	Il modello ordinario si recupera imponendo
	\[
	P_{m-2,h}=\cdots=P_{0,h}=I_b
	\]
	e
	\[
	\varepsilon_{m-1,h}
	=
	\cdots
	=
	\varepsilon_{0,h}.
	\]

	Per esponenti non uniformi \(J_h\) non commuta in generale con
	\(\MSC_m\). La forma fondamentale resta però utilizzabile direttamente:
	\[
	\begin{aligned}
	T_h
	&=
	P_h\circ\MSC_m\circ J_h\\
	&=
	P_h\circ\vartheta(J_h)\circ\MSC_m.
	\end{aligned}
	\]
	La rotazione degli esponenti è
	\[
	\vartheta(J_h)
	=
	R_b^{\varepsilon_{0,h}}
	\otimes
	R_b^{\varepsilon_{m-1,h}}
	\otimes\cdots\otimes
	R_b^{\varepsilon_{1,h}}.
	\]

	Prima di introdurre una notazione abbreviata, componiamo direttamente gli
	\(m\) stadi:
	\[
	\begin{aligned}
	T
	&=
	(P_{m-1}\circ\MSC_m\circ J_{m-1})
	\circ\cdots\circ
	(P_0\circ\MSC_m\circ J_0)\\
	&=
	\bigl(P_{m-1}\circ\vartheta(J_{m-1})\bigr)
	\circ
	\vartheta\!\bigl(P_{m-2}\circ\vartheta(J_{m-2})\bigr)
	\circ \cdots \circ\vartheta^{m-1}\!\bigl(P_0\circ\vartheta(J_0)\bigr),
	\end{aligned}
	\]
	dove la copia finale di \(\MSC_m^m\) è scomparsa perché
	\(\MSC_m^m=I\).

	A questo punto introduciamo il \emph{blocco effettivo}
	\[
	\boxed{
	A_h
	=
	P_h\circ\vartheta(J_h).
	}
	\]
	La forma dello stadio diventa
	\[
	\boxed{
	T_h=A_h\circ\MSC_m,
	}
	\qquad
	A_h\in\mathcal H_{b,m},
	\]
	mentre
	\[
	T_h\in\mathcal H_{b,m}\MSC_m.
	\]
	La composizione completa assume quindi la forma
	\[
	\boxed{
	T
	=
	A_{m-1}
	\circ\vartheta(A_{m-2})
	\circ\cdots\circ
	\vartheta^{m-1}(A_0).
	}
	\]

	\begin{proposition}[Uniformità del singolo stadio]
		Fissata una parola
		\[
		\boldsymbol{\varepsilon}_h
		=
		(\varepsilon_{m-1,h},\dots,\varepsilon_{0,h})
		\in\{0,1\}^m,
		\]
		al variare di \(P_h\in\mathcal H_{b,m}\) il blocco effettivo
		\[
		A_h=P_h\circ\vartheta(J_h)
		\]
		percorre una volta ciascun elemento di
		\(\mathcal H_{b,m}\). Al variare anche della parola binaria, ogni
		elemento del gruppo possiede esattamente \(2^m\) rappresentazioni
		mediante i controlli di uno stadio.
	\end{proposition}

	\begin{proof}
		Per la parola binaria fissata, \(\vartheta(J_h)\) è un elemento fissato
		di \(\mathcal H_{b,m}\). La moltiplicazione a destra
		\[
		P_h
		\longmapsto
		A_h=P_h\circ\vartheta(J_h)
		\]
		è quindi una biiezione del gruppo in sé. Le parole binarie possibili
		sono \(2^m\), e per ciascuna di esse il controllo \(P_h\) è determinato
		univocamente dall'operatore risultante.
	\end{proof}

	Diremo che una famiglia di controlli \emph{satura}
	\(\mathcal H_{b,m}\) quando l'insieme delle trasformazioni finali prodotte
	coincide con l'intero gruppo.

	\begin{theorem}[Saturazione e uniformità delle molteplicità del ciclo esteso]\label{thm:saturazione-uniformita-estesa}
		La famiglia tensoriale estesa soddisfa entrambe le proprietà seguenti.
		\begin{enumerate}
			\item L'insieme delle trasformazioni finali distinte è esattamente
			\(\mathcal H_{b,m}\), e quindi
			\[
			|\mathcal H_{b,m}|=(b!)^m.
			\]

			\item Ogni trasformazione finale \(T\in\mathcal H_{b,m}\) possiede
			esattamente
			\[
			|\mathcal H_{b,m}|^{m-1}
			=
			(b!)^{m(m-1)}
			\]
			decomposizioni
			\[
			T
			=
			A_{m-1}\circ\vartheta(A_{m-2})
			\circ\cdots\circ\vartheta^{m-1}(A_0).
			\]
			Si contano come uguali due decomposizioni quando coincidono, stadio per
			stadio, i valori dei blocchi effettivi \(A_h\). Per ogni decomposizione,
			i controlli locali possiedono \(2^{m^2}\) rappresentazioni complete.
			Pertanto ogni trasformazione finale possiede
			\[
			|\mathcal H_{b,m}|^{m-1}2^{m^2}
			\]
			storie complete dei controlli.
		\end{enumerate}
	\end{theorem}

	\begin{proof}
		Per la prima affermazione, ciascun blocco effettivo \(A_h\) può percorrere
		liberamente \(\mathcal H_{b,m}\), anche se le parole di rovesciamento vengono
		fissate in anticipo. Poiché \(\vartheta\) è un automorfismo del gruppo,
		ogni trasformazione prodotta appartiene a \(\mathcal H_{b,m}\).

		Reciprocamente, sia \(K\in\mathcal H_{b,m}\). Sono ammesse le scelte
		\[
		A_{m-1}=K,
		\qquad
		A_0=\cdots=A_{m-2}=I,
		\]
		perché, per ogni stadio, la mappa
		\(P_h\mapsto A_h=P_h\circ\vartheta(J_h)\) è biiettiva. La formula del
		ciclo si riduce allora a \(T=K\). Ogni elemento del gruppo è quindi
		raggiunto.

		Per la seconda affermazione, si scelgano liberamente i blocchi
		\(A_0,\dots,A_{m-2}\). Poiché \(\vartheta\) è un automorfismo di
		\(\mathcal H_{b,m}\), il blocco \(A_{m-1}\) è determinato univocamente
		dalla trasformazione finale assegnata. Le decomposizioni sono dunque
		\(|\mathcal H_{b,m}|^{m-1}\).

		Per ciascun blocco effettivo \(A_h\), la proposizione precedente fornisce
		\(2^m\) rappresentazioni mediante i controlli
		\((P_h,\boldsymbol{\varepsilon}_h)\). Gli \(m\) stadi sono indipendenti,
		per cui il fattore complessivo è
		\[
		(2^m)^m=2^{m^2}.
		\]
	\end{proof}

	\subsection{Formula locale del blocco effettivo}

	Nella famiglia estesa,
	\[
	A_h
	=
	(P_{m-1,h}\circ R_b^{\varepsilon_{0,h}})
	\otimes
	(P_{m-2,h}\circ R_b^{\varepsilon_{m-1,h}})
	\otimes\cdots\otimes
	(P_{0,h}\circ R_b^{\varepsilon_{1,h}}).
	\]
	Per evitare ambiguità nei casi piccoli, indicando con \(Q_{i,h}\) il
	fattore locale del blocco effettivo si può scrivere
	\[
	\boxed{
	Q_{i,h}
	=
	P_{i,h}\circ
	R_b^{\varepsilon_{(i+1)\bmod m,h}}.
	}
	\]
	Nel modello ordinario la stessa formula si specializza in
	\[
	A_h
	=
	(S_h\circ R_b^{\varepsilon_h})
	\otimes
	(R_b^{\varepsilon_h})^{\otimes(m-1)}.
	\]

	Per \(b=m=3\), ponendo \(R=R_3\),
	\begin{equation}\label{eq:Ah-ternario}
	A_h
	=
	(P_{2,h}\circ R^{\varepsilon_{0,h}})
	\otimes
	(P_{1,h}\circ R^{\varepsilon_{2,h}})
	\otimes
	(P_{0,h}\circ R^{\varepsilon_{1,h}}).
	\end{equation}
	Gli indici incrociati sono l'effetto diretto del trasporto \(\vartheta\).

\begin{remark}[Struttura delle composizioni parziali]
		La relazione \(\MSC_m^m=I\) mostra che l'ordine di \(\MSC_m\) divide
		\(m\). Per \(0<k<m\), la rotazione \(\MSC_m^k\) sposta il tensore
		con una sola cifra uguale a \(1\) e tutte le altre uguali a \(0\);
		quindi non è l'identità e l'ordine è esattamente \(m\).
		Inoltre \(\vartheta\) ruota i fattori di \(\mathcal H_{b,m}\), che è
		perciò stabile per coniugazione mediante \(\MSC_m\). Una potenza
		\(\MSC_m^k\), con \(0<k<m\), fa dipendere una cifra di uscita da un
		posto diverso dal corrispondente posto d'ingresso: variando quel posto e
		fissando la cifra d'ingresso si esclude la separabilità digitale. Pertanto
		\(\mathcal H_{b,m}\cap\langle\MSC_m\rangle=\{I\}\).
		Stabilità e intersezione banale danno la decomposizione del gruppo generato
		come prodotto semidiretto interno, e dunque
		\[
		\langle\mathcal H_{b,m},\MSC_m\rangle
		\cong\mathcal H_{b,m}\rtimes C_m,
		\]
		dove \(C_m\) agisce ruotando ciclicamente i fattori. Per
		\(1\le k\le m\),
		\[
				T_{k-1}\circ\cdots\circ T_0
		\in
		\mathcal H_{b,m}\MSC_m^k.
		\]
		In particolare, dopo \(m\) stadi si ritorna in
		\(\mathcal H_{b,m}\).
	\end{remark}

\section{Ricostruzione da statistiche aggregate sulle fibre}

	La caratterizzazione precedente ricostruisce le componenti locali leggendo le
	cifre delle immagini. Qui consideriamo invece le somme di fibra come una
	rappresentazione aggregata della permutazione. La progressione è dalla
	formula per gli elementi separabili alla ricostruzione entro quella classe,
	fino al riconoscimento fra tutte le permutazioni di \(X=\{0,\dots,b^m-1\}\).

	La ricostruzione di matrici di permutazione da dati aggregati compare anche
	nello studio degli X-rays, definiti mediante somme antidiagonali
	\cite{bebeacua-et-al}. In quel contesto permutazioni diverse possono avere
	lo stesso X-ray, e l'unicità richiede un'analisi specifica. Le somme sulle
	fibre digitali utilizzate qui sono statistiche diverse: il criterio seguente
	riguarda queste partizioni e la classe separabile, senza risolvere un problema
	generale di tomografia discreta (\emph{discrete tomography}).

	Per una permutazione digitalmente separabile
	\(\pi=\pi_{\boldsymbol{\rho}}\), si definisce
	\[
	\Sigma_{i,t}(\pi)
	=
	\sum_{x\in F_{i,t}}\pi(x).
	\]
	La somma raccoglie le immagini di tutte le posizioni che condividono la cifra
	\(t\) al livello \(i\).

	\begin{theorem}[Formula delle statistiche di fibra]\label{thm:statistiche-fibra}
		Per ogni \(i\in\{0,\dots,m-1\}\) e \(t\in\Db\),
		\[
		\Sigma_{i,t}(\pi)
		=
		C_i+b^{m-1+i}\rho_i(t),
		\]
		dove
		\[
		C_i
		=
		b^{m-2}\frac{b(b-1)}2
		\sum_{\substack{0\le j\le m-1\\j\ne i}}b^j.
		\]
	\end{theorem}

	\begin{proof}
		Se \(x=\val(a_{m-1},\dots,a_0)\), allora
		\[
		\pi(x)
		=
		\sum_{j=0}^{m-1}b^j\rho_j(a_j).
		\]
		Sulla fibra \(F_{i,t}\), la cifra \(a_i\) è fissata. Il contributo del
		livello \(i\) alla somma è quindi
		\[
		\sum_{x\in F_{i,t}}b^i\rho_i(t)
		=
		b^{m-1+i}\rho_i(t).
		\]
		Per ogni \(j\ne i\), la cifra \(a_j\) assume ciascun valore di \(\Db\)
		esattamente \(b^{m-2}\) volte. Poiché \(\rho_j\) è una permutazione,
		\[
		\sum_{u=0}^{b-1}\rho_j(u)
		=
		\sum_{u=0}^{b-1}u
		=
		\frac{b(b-1)}2.
		\]
		Sommando i contributi dei livelli \(j\ne i\) si ottiene \(C_i\), che non
		dipende da \(t\) né dalle permutazioni locali diverse da \(\rho_i\).
	\end{proof}

	La formula può essere riscritta in termini di media sulla fibra:
	\[
	\boxed{
	\frac{\Sigma_{i,t}(\pi)}{b^{m-1}}
	=
	\frac{b^m-1}{2}
	+b^i\left(\rho_i(t)-\frac{b-1}{2}\right).
	}
	\]
	Questa forma separa il valore medio globale dal contributo del livello
	digitale \(i\).

	\begin{corollary}[Sufficienza delle statistiche aggregate]
		All'interno del gruppo delle permutazioni digitalmente separabili, la
		famiglia completa delle quantità \(\Sigma_{i,t}\) determina univocamente la
		permutazione globale. Le componenti locali sono ricostruite mediante
		\[
		\rho_i(t)
		=
		\frac{\Sigma_{i,t}(\pi)-C_i}{b^{m-1+i}}.
		\]
		In particolare, la mappa
		\[
		\pi
		\longmapsto
		\bigl(\Sigma_{i,t}(\pi)\bigr)_{
			0\le i\le m-1,\;t\in\Db}
		\]
		è iniettiva sulla classe digitalmente separabile.
	\end{corollary}

	\begin{proof}
		La formula ricostruisce tutti i valori di tutte le permutazioni
		\(\rho_i\). Il teorema di caratterizzazione determina quindi
		\(\pi=\pi_{\boldsymbol{\rho}}\) in modo unico.
	\end{proof}

	\begin{remark}
		La famiglia completa delle statistiche \(\Sigma_{i,t}\) è sufficiente, ma
		non se ne afferma la minimalità. Per ciascun livello, una volta ricostruiti
		\(b-1\) valori della permutazione locale, l'ultimo è determinato dai
		precedenti.
	\end{remark}

	Il risultato seguente supera la ricostruzione interna del Corollario~6.2:
	se una permutazione arbitraria di \(X\) produce somme di fibra compatibili
	con un elemento separabile, deve coincidere con esso. L'unicità vale quindi
	nell'intero \(S_{b^m}\), senza assumere la separabilità della permutazione
	da riconoscere.

	\begin{theorem}[Criterio di riconoscimento mediante somme di fibra]\label{thm:criterio-fibra}
		Sia \(\pi\in S_{b^m}\) una permutazione arbitraria. Per ogni \(i,t\)
		calcoliamo
		\[
		\Sigma_{i,t}(\pi)=\sum_{x\in F_{i,t}}\pi(x),
		\qquad
		\rho_i(t)=\frac{\Sigma_{i,t}(\pi)-C_i}{b^{m-1+i}}.
		\]
		Chiamiamo \(\rho_i(t)\) i \emph{valori candidati} ricostruiti dalle
		statistiche. Allora
		\[
		\boxed{
		\pi\ \text{è digitalmente separabile}
		\iff
		\bigl(\rho_i(t)\bigr)_{t\in D_b}
		\ \text{è una permutazione di }D_b
		\ \text{per ogni }i.
		}
		\]
		In tal caso
		\[
		\pi(x)=\sum_{i=0}^{m-1}b^i\rho_i(d_i(x)).
		\]
		Equivalentemente, nessuna permutazione diversa può avere le stesse
		somme di fibra di una permutazione digitalmente separabile.
	\end{theorem}

	\begin{proof}
		La necessità segue dal Teorema~\ref{thm:statistiche-fibra}. Per la
		sufficienza, supponiamo che i valori candidati ricostruiti definiscano, per
		ogni livello \(i\), una permutazione \(\rho_i\) di \(D_b\). La funzione
		\[
		f(x)=\sum_{i=0}^{m-1}b^i\rho_i(d_i(x))
		\]
		è allora una permutazione digitalmente separabile. Dalla formula delle
		statistiche di fibra e dalla definizione di \(\rho_i(t)\) segue
		\[
		\Sigma_{i,t}(f)
		=
		C_i+b^{m-1+i}\rho_i(t)
		=
		\Sigma_{i,t}(\pi).
		\]
		Quindi l'uguaglianza delle somme di fibra dà
		\[
		\begin{aligned}
		\sum_x\pi(x)f(x)
		&=\sum_i b^i\sum_t\rho_i(t)\Sigma_{i,t}(\pi)\\
		&=\sum_i b^i\sum_t\rho_i(t)\Sigma_{i,t}(f)\\
		&=\sum_x f(x)^2.
		\end{aligned}
		\]
		Poiché \(\pi\) e \(f\) permutano lo stesso insieme
		\(\{0,\dots,b^m-1\}\),
		\[
		\sum_x\pi(x)^2=\sum_x f(x)^2.
		\]
		Pertanto
		\[
		\sum_x\bigl(\pi(x)-f(x)\bigr)^2=0,
		\]
		da cui \(\pi=f\).
	\end{proof}

\paragraph*{Unicità additiva e rappresentazione per altezze.}
Il principio di unicità del Teorema~\ref{thm:criterio-fibra} può essere
collegato al Lemma~2.14 di Kemperman \cite[p.~282]{kemperman}.
Poniamo \(N=b^m\) e identifichiamo le posizioni con gli indirizzi
\(a=(a_{m-1},\ldots,a_0)\in\Db^m\). In questa rappresentazione leggiamo
ogni permutazione \(g\) come funzione degli indirizzi, scrivendo \(g(a)\)
per il valore \(g(\val(a))\).

Nell'insieme finito
\[
\Omega
=
\Db^m\times\{0,\ldots,N-2\},
\]
consideriamo la rappresentazione per altezze
\[
S_g
=
\{(a,h)\in\Omega:\ h<g(a)\}.
\]
Il marginale di conteggio relativo alla coordinata \(a_i\), nel valore
\(u\), è \(\Sigma_{i,u}(g)\), ossia la somma dei valori di \(g\) sulla
fibra \(F_{i,u}\). Il marginale dell'altezza, nel valore \(h\), è
\(N-1-h\), poiché \(g\) permuta i valori \(0,\ldots,N-1\).

Se
\[
f(a)=\sum_{i=0}^{m-1}b^i\rho_i(a_i),
\]
allora
\[
S_f
=
\left\{
(a,h)\in\Omega:
\sum_{i=0}^{m-1}b^i\rho_i(a_i)-h-\frac12\ge0
\right\}.
\]
La funzione che definisce la soglia è una somma di funzioni delle singole
coordinate; nell'ambiente finito con misura di conteggio, \(S_f\) soddisfa
quindi le ipotesi di additività del lemma.

Il risultato di Kemperman rende \(S_f\) unico fra gli insiemi con gli stessi
marginali, senza richiedere l'additività del concorrente. Pertanto, se \(g\)
e \(f\) hanno le stesse somme di fibra, anche i marginali dell'altezza
coincidono, da cui \(S_g=S_f\) e infine \(g=f\).
Questa rappresentazione esplicita il rapporto con il principio generale
di unicità; la dimostrazione quadratica precedente fornisce direttamente
la conclusione nel modello digitale.

	\begin{remark}[Ricostruzione e perdita della storia dei controlli]
		Le statistiche di fibra ricostruiscono univocamente il codice locale
		effettivo e quindi la trasformazione finale, anche quando essa proviene
		dalla famiglia tensoriale estesa. Non determinano invece una particolare
		decomposizione nei blocchi effettivi né una specifica storia dei controlli:
		per il Teorema~\ref{thm:saturazione-uniformita-estesa}, ogni trasformazione
		finale possiede \(|\mathcal H_{b,m}|^{m-1}\) decomposizioni nei blocchi
		effettivi e ciascuna decomposizione ammette \(2^{m^2}\) storie complete dei
		controlli.
	\end{remark}

	\begin{example}[Un ciclo e la ricostruzione in base due]
		Per \(b=2\), \(m=3\), scegliamo, senza rovesciamenti,
		\(S_0=(0\,1)\), \(S_1=I_2\), \(S_2=(0\,1)\). Il ciclo realizza
		\[
		\pi\bigl((a_2a_1a_0)_2\bigr)=\val(1-a_2,a_1,1-a_0),
		\qquad P_\pi=R_2\otimes I_2\otimes R_2.
		\]
		La posizione \(2=(010)_2\) segue \(2\mapsto5\mapsto6\mapsto7\).
		La lista delle immagini, nell'ordine \(x=0,\dots,7\), è
		\((5,4,7,6,1,0,3,2)\). Le coppie di somme
		\((\Sigma_{i,0},\Sigma_{i,1})\), per \(i=0,1,2\), sono
		\((16,12)\), \((10,18)\), \((22,6)\). Con
		\((C_0,C_1,C_2)=(12,10,6)\) e coefficienti \((4,8,16)\), la
		formula restituisce \(\rho_0=(0\,1)\), \(\rho_1=\mathrm{id}\),
		\(\rho_2=(0\,1)\); il Teorema~\ref{thm:criterio-fibra} certifica
		anche l'unicità di questa permutazione in \(S_8\).
	\end{example}

	\begin{example}[Rifiuto di una permutazione non separabile]
		Ancora per \(b=2\), \(m=3\), sia \(\pi=(0\,1)\in S_8\), con
		le altre posizioni fisse. Sulla fibra \(F_{0,0}=\{0,2,4,6\}\), le
		immagini \(1\) e \(2\) hanno cifre meno significative diverse, perciò
		\(\pi\) non è digitalmente separabile. Le somme al livello \(0\) sono
		\[
		\Sigma_{0,0}=1+2+4+6=13,\qquad
		\Sigma_{0,1}=0+3+5+7=15.
		\]
		Poiché \(C_0=12\) e \(b^{m-1}=4\), i candidati sono
		\(\rho_0(0)=1/4\) e \(\rho_0(1)=3/4\). Non formano una
		permutazione di \(D_2\), e il test rifiuta \(\pi\) già a questo livello.
	\end{example}

	\section{Il caso delle ventisette carte}

	Per \(b=m=3\), le posizioni sono numerate da \(0\) a \(26\) e ogni posizione
	ha la forma
	\[
	x=(a_2a_1a_0)_3.
	\]
	Le fibre digitali contengono nove posizioni. La formula generale diventa
	\[
	\Sigma_{i,t}(\pi)
	=
	C_i+3^{2+i}\rho_i(t),
	\]
	con
	\[
	C_0=108,
	\qquad
	C_1=90,
	\qquad
	C_2=36.
	\]
	I valori possibili sono
	\[
	\begin{array}{c@{\qquad}ccc}
		\toprule
		i & \rho_i(t)=0 & \rho_i(t)=1 & \rho_i(t)=2\\
		\midrule
		0 & 108 & 117 & 126\\
		1 &  90 & 117 & 144\\
		2 &  36 & 117 & 198\\
		\bottomrule
	\end{array}
	\]
	Le differenze consecutive sono \(9\), \(27\) e \(81\), in accordo con i
	coefficienti \(3^{2+i}\).

	\begin{example}[Ricostruzione di un codice]
		Si consideri il codice
		\[
		\boldsymbol{\rho}
		=
		(\rho_2,\rho_1,\rho_0),
		\]
		con
		\[
		\rho_0=(0\,1\,2),
		\qquad
		\rho_1=(0\,2),
		\qquad
		\rho_2=(0\,1),
		\]
		dove i cicli agiscono su \(\{0,1,2\}\). Le statistiche di fibra sono
		\[
		\begin{array}{c@{\qquad}ccc}
			\toprule
			i & \Sigma_{i,0} & \Sigma_{i,1} & \Sigma_{i,2}\\
			\midrule
			0 & 117 & 126 & 108\\
			1 & 144 & 117 &  90\\
			2 & 117 &  36 & 198\\
			\bottomrule
		\end{array}
		\]
		La formula di ricostruzione restituisce, per il livello \(0\),
		\[
		(\rho_0(0),\rho_0(1),\rho_0(2))
		=
		\left(
		\frac{117-108}{9},
		\frac{126-108}{9},
		\frac{108-108}{9}
		\right)
		=
		(1,2,0).
		\]
		Analogamente,
		\[
		(\rho_1(0),\rho_1(1),\rho_1(2))=(2,1,0),
		\qquad
		(\rho_2(0),\rho_2(1),\rho_2(2))=(1,0,2).
		\]
		Le tre permutazioni locali effettive vengono quindi ricostruite
		integralmente dai nove dati aggregati.
	\end{example}

	Nel caso \(b=m=3\), ponendo
	\[
	R=
	\begin{pmatrix}
	0&0&1\\
	0&1&0\\
	1&0&0
	\end{pmatrix},
	\]
	il rovesciamento globale ordinario è il sottocaso diagonale
	\[
	J_h
	=
	R^{\varepsilon_h}\otimes
	R^{\varepsilon_h}\otimes
	R^{\varepsilon_h}.
	\]
	Con le regole ordinarie del gioco, a ogni stadio si sceglie una sola
	raccolta fra le sei possibili e una sola scelta binaria di rovesciamento globale. Le
	procedure complete sono quindi
	\[
	(6\cdot2)^3=1728,
	\]
	mentre le trasformazioni globali distinte restano
	\[
	6^3=216.
	\]
	Le otto procedure associate a una stessa trasformazione finale emergono
	come molteplicità completa del ciclo: il fattore \(8=2^3\) appartiene alla
	storia globale del procedimento e non a una molteplicità locale del singolo
	stadio. Ciò non esclude molteplicità già presenti in prefissi del ciclo.

	\begin{example}[Una traiettoria ordinaria completa]
		Prendiamo
		\[
		S_0=(0\,1\,2),\qquad S_1=(0\,1),\qquad S_2=(0\,2),
		\]
		e la storia
		\[
		(\varepsilon_0,\varepsilon_1,\varepsilon_2)=(1,0,1).
		\]
		La posizione \(5=(012)_3\) segue
		\[
		5\longmapsto16\longmapsto5\longmapsto25.
		\]
		Il Corollario~\ref{cor:formula-inversione-raccolte} fornisce
		\[
		Q_0=(0\,2\,1),\qquad Q_1=(1\,2),\qquad Q_2=(0\,2),
		\]
		e quindi
		\[
		(012)_3\longmapsto(221)_3=25.
		\]
	\end{example}

	Nella famiglia tensoriale estesa, in ciascuno stadio si scelgono
	indipendentemente tre permutazioni locali e tre fattori binari di
	rovesciamento; il blocco effettivo è quello della
	formula~\eqref{eq:Ah-ternario}. I conteggi essenziali sono:
	\[
	\begin{array}{l@{\qquad}r}
	\toprule
	\text{configurazioni per stadio} & 6^3 2^3=1728\\
	\text{blocchi effettivi distinti} & 216\\
	\text{rappresentazioni per blocco} & 8\\
	\text{storie complete} & 1728^3=5\,159\,780\,352\\
	\text{decomposizioni per trasformazione finale} & 216^2=46\,656\\
	\text{storie per decomposizione} & 8^3=512\\
	\text{storie per trasformazione finale} & 23\,887\,872\\
	\bottomrule
	\end{array}
	\]
	Le trasformazioni finali distinte restano \(216\). L'otto-a-uno del modello
	ordinario è una molteplicità del ciclo completo; nella famiglia estesa il
	fattore \(8\) compare già nella rappresentazione di ciascun blocco effettivo.

	Nel linguaggio del gioco, la tabella delle immagini delle ventisette
	posizioni rappresenta la trasformazione globale. Nel sottocaso ordinario senza rovesciamenti, il codice
	di raccolta e il codice locale effettivo coincidono e la corrispondenza con la
	tabella delle immagini è biunivoca. Quando si registra anche il rovesciamento globale, la stessa
	trasformazione finale ammette invece otto storie fisiche: le statistiche di
	fibra ricostruiscono la trasformazione finale e i fattori effettivi, non la particolare
	storia di rovesciamenti che li ha prodotti. Se tale storia è nota, le raccolte
	fisiche si ricostruiscono univocamente dal Corollario~\ref{cor:formula-inversione-raccolte}.

	\section{Conseguenze strutturali}

		La fattorizzazione separata rende immediatamente leggibili alcuni invarianti
	della trasformazione globale. Indichiamo con \(\fix(\sigma)\) il numero dei
	punti fissi della permutazione \(\sigma\).

	\begin{proposition}
		Per
		\[
		\boldsymbol{\rho}
		=(\rho_{m-1},\dots,\rho_0)\in S_b^m
		\]
		valgono
		\[
		\ord(\pi_{\boldsymbol{\rho}})
		=
		\mcm\bigl(
		\ord(\rho_{m-1}),\dots,\ord(\rho_0)
		\bigr)
		\]
		e
		\[
		\fix(\pi_{\boldsymbol{\rho}})
		=
		\prod_{i=0}^{m-1}\fix(\rho_i).
		\]
	\end{proposition}

	\begin{proof}
		La potenza \(\pi_{\boldsymbol{\rho}}^n\) è l'identità esattamente quando
		\(\rho_i^n\) è l'identità per ogni \(i\); il minimo intero positivo con
		questa proprietà è il minimo comune multiplo degli ordini locali. Una
		posizione è fissa quando ciascuna delle sue cifre è fissata dalla
		permutazione del livello corrispondente. Le scelte delle cifre sono
		indipendenti, e le cardinalità si moltiplicano.
	\end{proof}

	Gli invarianti della proposizione dipendono quindi soltanto dai fattori
	locali effettivi della trasformazione completa. Due procedure fisiche che
	differiscono per la storia dei rovesciamenti ma producono lo stesso codice
	effettivo hanno necessariamente lo stesso ordine, gli stessi punti fissi e,
	più in generale, gli stessi invarianti della permutazione globale.

	\begin{remark}
		In \(S_b\) le classi di coniugio sono determinate dal tipo di ciclo e sono
		quindi indicizzate dalle partizioni di \(b\). Le classi di coniugio del
		gruppo interno \(S_b^m\) sono scelte indipendentemente nei fattori e sono
		perciò \(p(b)^m\), dove \(p(b)\) è il numero delle partizioni di \(b\). Per \(b=m=3\) se ne ottengono \(3^3=27\).
		Questa classificazione riguarda il gruppo digitalmente separabile; la
		coniugazione nel gruppo ambiente \(S_{b^m}\) può fondere classi interne
		distinte di \(S_b^m\).
	\end{remark}

	\section{Discussione conclusiva}

	\subsection{Risultati e portata del modello}

	Il modello non adattivo su \(N=b^m\) posizioni descrive ogni stadio di
	distribuzione e raccolta come una rotazione delle coordinate digitali
	seguita da un'azione locale. Dopo \(m\) stadi la rotazione si chiude e gli
	esiti costituiscono precisamente il gruppo delle trasformazioni
	digitalmente separabili, isomorfo a \(S_b^m\). La preservazione delle
	partizioni digitali caratterizza questa classe e determina univocamente
	le sue componenti locali.

	Le somme di fibra sviluppano il problema inverso. Esse consentono di
	recuperare i fattori di una trasformazione separabile e, attraverso il
	Teorema~\ref{thm:criterio-fibra}, di riconoscere tale struttura fra tutte le
	permutazioni delle \(b^m\) posizioni. L'unicità riguarda la trasformazione
	finale e il suo codice locale effettivo; le diverse storie operative che
	realizzano lo stesso esito conservano le molteplicità determinate nella
	sezione~5.

	\subsection{Rapporto con la letteratura e formulazione digitale}

	Gli strumenti della trattazione appartengono a quadri matematici
	precedenti. Il rapporto fra giochi di Gergonne, notazione posizionale e
	ordinamento per cifre è discusso da Bolker \cite{bolker}. La teoria dei
	mescolamenti perfetti è sviluppata da Diaconis, Graham e Kantor
	\cite{diaconis-graham-kantor}. Matrici di permutazione, prodotti di
	Kronecker, matrici di commutazione e riordinamenti dei fattori
	costituiscono il linguaggio algebrico utilizzato qui; per questi ultimi
	strumenti si rinvia a Henderson e Searle \cite{henderson-searle} e a
	D'Angeli e Donno \cite{dangeli-donno}. L'azione indipendente sulle
	coordinate realizza il gruppo di base di un prodotto intrecciato in
	azione prodotto, nel quadro descritto da Praeger e Schneider
	\cite{praeger-schneider}.

	Anche la determinazione di insiemi discreti mediante le proiezioni sugli
	assi appartiene alla letteratura precedente, nella quale si colloca il
	lavoro di Fishburn, Lagarias, Reeds e Shepp \cite{fishburn-et-al}.
	Per il principio di unicità pertinente al presente modello, il riferimento
	utilizzato è il Lemma~2.14 di Kemperman \cite[p.~282]{kemperman}:
	un insieme additivo è determinato dai propri marginali. L'osservazione
	successiva al Teorema~\ref{thm:criterio-fibra} esplicita il collegamento,
	rappresentando i valori della permutazione mediante altezze.

	Nel modello digitale considerato, la separabilità conduce alla formula
	\[
	\Sigma_{i,u}
	=
	C_i+b^{m-1+i}\rho_i(u),
	\]
	con le costanti \(C_i\) determinate nella sezione~6. L'inversione affine
	\[
	\rho_i(u)
	=
	\frac{\Sigma_{i,u}-C_i}{b^{m-1+i}}
	\]
	recupera le componenti locali. Per una permutazione arbitraria, il
	controllo che ogni lista candidata sia una permutazione di \(\Db\)
	costituisce il criterio operativo di riconoscimento: quando il controllo
	è soddisfatto, i candidati ricostruiscono la permutazione assegnata.

	La dimostrazione quadratica del Teorema~\ref{thm:criterio-fibra} rimane
	autonoma. L'argomento può anche essere letto geometricamente:
	l'uguaglianza delle somme di fibra rende il residuo \(\pi-f\) ortogonale
	al candidato separabile \(f\), mentre il comune multinsieme dei valori
	fissa \(\|\pi\|=\|f\|\). L'identità pitagorica impone allora \(\pi=f\).

	La struttura \(S_b^m\) e il principio generale di unicità hanno dunque
	precedenti. La presente trattazione ne sviluppa una formulazione digitale
	esplicita, collegando le somme di fibra, le costanti \(C_i\), il recupero
	delle componenti locali e il controllo di riconoscimento al procedimento
	di distribuzione e raccolta studiato. Nelle fonti direttamente consultate
	non è stata individuata una formulazione operativa completa che riunisca
	questi elementi. Questa constatazione delimita il confronto bibliografico
	effettuato; non viene qui avanzata una rivendicazione di novità o di
	priorità della procedura.

	\subsection{Limiti e sviluppi}

	I risultati dipendono dalle operazioni ammesse e dalla struttura digitale
	fissata. I rovesciamenti globali e l'estensione tensoriale con controlli
	locali indipendenti conservano la classe degli esiti finali,
	modificandone le molteplicità di realizzazione. L'estensione tensoriale
	ha statuto algebrico e non presuppone che ogni controllo sia realizzabile
	mediante le normali raccolte.

	La disposizione rettangolare mediante due fattori di dimensioni
	differenti richiede invece un'analisi distinta. Lo scambio dei fattori
	descrive un passaggio invertibile, ma da solo non garantisce la chiusura
	ciclica del modello omogeneo. Nel gioco delle ventuno carte occorre
	inoltre distinguere la permutazione fisica del mazzo dalla contrazione
	delle possibilità relative alla carta scelta.

	La sufficienza delle statistiche di fibra non ne dimostra la minimalità
	né la stabilità rispetto a dati perturbati. Lo studio di questi aspetti,
	di altre famiglie di partizioni e di ulteriori classi riconoscibili
	mediante dati aggregati costituisce uno sviluppo successivo della
	trattazione.

\appendix

	\section{Il caso rettangolare
		\texorpdfstring{\(N=i\,j\)}{N=ij} e la dinamica informativa}

	La fattorizzazione \(N=b^m\) utilizza livelli omogenei. Nei procedimenti
	rettangolari il numero delle carte assume invece la forma
	\[
	N=i\,j,
	\qquad i,j\in\mathbb N_{>0},
	\]
	con fattori che possono essere differenti. Se il mazzo viene distribuito
	ciclicamente in \(i\) colonne di lunghezza \(j\), ogni posizione si scrive
	\[
	x=iq+r,
	\qquad
	0\le r<i,
	\qquad
	0\le q<j.
	\]
	Il resto \(r\) identifica la colonna e il quoziente \(q\) la quota al suo
	interno. Alla posizione si associa
	\[
	e_q^{(j)}\otimes e_r^{(i)}
	\in
	\mathbb R^j\otimes\mathbb R^i.
	\]

	Lo scambio fra quota e colonna è realizzato dalla matrice di commutazione
	\[
	K_{j,i}:
	\mathbb R^j\otimes\mathbb R^i
	\longrightarrow
	\mathbb R^i\otimes\mathbb R^j,
	\]
	definita da
	\[
	K_{j,i}
	\bigl(e_q^{(j)}\otimes e_r^{(i)}\bigr)
	=
	e_r^{(i)}\otimes e_q^{(j)}.
	\]
	Vale
	\[
	K_{j,i}^{-1}=K_{i,j}=K_{j,i}^{\mathsf T}.
	\]

	Sia \(\tau\in S_i\) la permutazione che assegna alla colonna \(r\) la
	posizione di blocco \(\tau(r)\), e sia \(P_\tau\) la sua matrice. Conservando
	l'ordine interno delle colonne, il passaggio fisico è
	\[
	\pi_\tau(iq+r)
	=
	j\tau(r)+q
	\]
	e il corrispondente operatore è
	\[
	M_\tau
	=
	(P_\tau\otimes I_j)\circ K_{j,i}:
	\mathbb R^j\otimes\mathbb R^i
	\longrightarrow
	\mathbb R^i\otimes\mathbb R^j.
	\]
	Dopo aver identificato i due prodotti tensoriali ordinati con
	\(\mathbb R^{ij}\) mediante le rispettive basi linearizzate, \(M_\tau\) è
	rappresentato da una matrice di permutazione \(ij\times ij\). La sua
	invertibilità descrive la reversibilità del passaggio fisico.

	L'operatore \(K_{j,i}\) e la rotazione \(\MSC_m\) esprimono entrambi una
	permutazione dei fattori. Nel caso \(b^m\), tuttavia, tutti i fattori hanno la
	stessa dimensione e una rotazione di ordine \(m\) espone successivamente ogni
	cifra della medesima base. Nel caso rettangolare, quota e colonna hanno in
	generale dimensioni differenti; lo scambio è invertibile, ma non produce da
	solo una chiusura digitale analoga.

	Nel gioco delle ventuno carte l'ordine di raccolta è inoltre adattivo. Dopo
	ogni distribuzione il partecipante indica la colonna contenente la carta
	scelta e quella colonna viene collocata al centro. Se, più in generale, la
	colonna indicata viene collocata nella posizione di blocco
	\(c\in\{0,\dots,i-1\}\), la posizione della sola carta scelta segue la mappa
	\[
	F_c(x)
	=
	jc+\left\lfloor\frac{x}{i}\right\rfloor.
	\]
	Questa non è una singola permutazione fisica del mazzo: per ogni posizione
	ipotetica seleziona il ramo corrispondente alla colonna che verrebbe indicata,
	e soddisfa
	\[
	F_c(iq)=F_c(iq+1)=\cdots=F_c(iq+i-1)=jc+q.
	\]
	Per \(i\geq2\) la mappa è quindi molti-a-uno. Nel caso degenere \(i=1\),
	invece, essa è biiettiva.

	Per il gioco classico delle ventuno carte,
	\[
	i=3,
	\qquad
	j=7,
	\qquad
	c=1,
	\]
	si ha
	\[
	F(x)=7+\left\lfloor\frac{x}{3}\right\rfloor.
	\]
	Partendo da
	\[
	X_0=\{0,1,\dots,20\},
	\]
	definiamo ricorsivamente \(X_{k+1}=F(X_k)\). Considerando in questo modo
	l'insieme complessivo delle posizioni raggiungibili al variare delle possibili
	storie di risposta, si ottiene
	\[
	X_1=\{7,8,\dots,13\},
	\qquad
	X_2=\{9,10,11\},
	\qquad
	X_3=\{10\},
	\]
	e quindi
	\[
	21\longmapsto7\longmapsto3\longmapsto1.
	\]
	Questi insiemi non vanno confusi con quelli condizionati a una specifica
	sequenza di risposte osservate, che possono essere più piccoli e si ottengono
	intersecando a ogni stadio con la colonna effettivamente indicata prima di
	applicare \(F\). La permutazione fisica del mazzo resta invertibile a ogni
	stadio; ciò che si contrae è l'incertezza relativa alla carta scelta. Questa distinzione separa
	la dinamica del mazzo dalla dinamica informativa e segna il confine naturale
	della struttura digitalmente separabile sviluppata nel corpo dell'articolo.

	\begin{thebibliography}{99}

		\bibitem{bolker}
		E.~D. Bolker,
		\emph{Gergonne's Card Trick, Positional Notation, and Radix Sort},
		Mathematics Magazine \textbf{83} (2010), no.~1, 46--49.
		\href{https://doi.org/10.4169/002557010X479983}{doi:10.4169/002557010X479983}.

		\bibitem{champanerkar-jani}
		J. Champanerkar e M. Jani,
		\emph{Using a Card Trick to Illustrate Fixed Points and Stability},
		PRIMUS \textbf{25} (2015), no.~5, 473--483.
		\href{https://doi.org/10.1080/10511970.2015.1031308}{doi:10.1080/10511970.2015.1031308}.

		\bibitem{diaconis-graham-kantor}
		P. Diaconis, R.~L. Graham e W.~M. Kantor,
		\emph{The Mathematics of Perfect Shuffles},
		Advances in Applied Mathematics \textbf{4} (1983), 175--196.
		\href{https://doi.org/10.1016/0196-8858(83)90009-X}{doi:10.1016/0196-8858(83)90009-X}.

		\bibitem{dangeli-donno}
		D. D'Angeli e A. Donno,
		\emph{Shuffling matrices, Kronecker product and Discrete Fourier Transform},
		Discrete Applied Mathematics \textbf{233} (2017), 1--18.
		\href{https://doi.org/10.1016/j.dam.2017.08.018}{doi:10.1016/j.dam.2017.08.018}.

		\bibitem{henderson-searle}
		H.~V. Henderson e S.~R. Searle,
		\emph{The Vec-Permutation Matrix, the Vec Operator and Kronecker Products: A Review},
		Linear and Multilinear Algebra \textbf{9} (1981), no.~4, 271--288.
		\href{https://doi.org/10.1080/03081088108817379}{doi:10.1080/03081088108817379}.

		\bibitem{praeger-schneider}
		C.~E. Praeger e C. Schneider,
		\emph{Embedding permutation groups into wreath products in product action},
		Journal of the Australian Mathematical Society \textbf{92} (2012),
		no.~1, 127--136.
		\href{https://doi.org/10.1017/S1446788712000110}{doi:10.1017/S1446788712000110}.

		\bibitem{bebeacua-et-al}
		C. Bebeacua, T. Mansour, A. Postnikov e S. Severini,
		\emph{On the X-rays of permutations},
		Electronic Notes in Discrete Mathematics \textbf{20} (2005), 193--203.
		\href{https://doi.org/10.1016/j.endm.2005.05.063}{doi:10.1016/j.endm.2005.05.063}.

		\bibitem{quintero}
		R. Quintero,
		\emph{On a Mathematical Model for an Old Card Trick},
		Recreational Mathematics Magazine \textbf{4} (2017), no.~7, 65--77.
		\href{https://doi.org/10.1515/rmm-2017-0014}{doi:10.1515/rmm-2017-0014}.


		\bibitem{kemperman}
		J.~H.~B. Kemperman,
		\emph{Sets of uniqueness and systems of inequalities having a unique solution},
		Pacific Journal of Mathematics \textbf{148} (1991), no.~2, 275--301.
		\href{https://doi.org/10.2140/pjm.1991.148.275}{doi:10.2140/pjm.1991.148.275}.

		\bibitem{fishburn-et-al}
		P.~C. Fishburn, J.~C. Lagarias, J.~A. Reeds e L.~A. Shepp,
		\emph{Sets uniquely determined by projections on axes II: Discrete case},
		Discrete Mathematics \textbf{91} (1991), no.~2, 149--159.
		\href{https://doi.org/10.1016/0012-365X(91)90106-C}{doi:10.1016/0012-365X(91)90106-C}.

	\end{thebibliography}
